\documentclass[11pt]{article}

\usepackage[margin=1.5in]{geometry}
\usepackage{lecnotes}
% \usepackage{mpass}
% \lstset{style=verb,language=mpass}
\input{lmacros}
% \usepackage[clock]{ifsym}
% \newcommand{\tm}{\Interval}
\newcommand{\tm}{t}

\newcommand{\course}{15-316: Software Foundations of Security \& Privacy}
\newcommand{\lecturer}{Matt Fredrikson}
\newcommand{\lecdate}{March 12, 2026}
\newcommand{\lecnum}{15}
\newcommand{\lectitle}{Timing Attacks}
\newcommand{\courseurl}{https://15316-cmu.github.io/2026/}

\begin{document}

\maketitle

\section{Introduction}

Deducing secrets (like passwords or private keys) by observing how long a
program runs is a commonly used attack on computer security.  This is true
despite the uncertainties about how long a program will actually run (which is
only stochastically connected to the source code).  For example,
\citet{Brumley05cn} demonstrate that remote timing attacks are quite feasible
and can be launched against security-critical code such as an implementation of
SSL\@.  They had to run a program many times, but in the end they were
successful in exploiting data-dependent optimizations in the cryptographic
primitives underlying SSL\@.  Due to the serious nature of such security
threats, defenses have been devised.  The primary ones are (1) introduce
randomness into the computation so that the time variations don't give away the
secret information, and (2) sharpen the information flow type system to ensure
``constant-time'' computation.  We explore the second and briefly touch on the
first.  Timing attacks against cryptographic primitives are still being
discovered even today.

There are other so-called \emph{side channels} for information.  Here are
some:
\begin{itemize}
\item Power consumption
\item Memory access patterns
\item Sequence of instructions executed
\item Electromagnetic radiation emitted from the hardware
\end{itemize}
Some of them require direct access to some hardware, other can be launched
remotely.  What they have in common in that they use physical properties of what
takes place in a computer when code is executed that is outside of the usual
abstract model of computation we work with when reasoning about our code.  This
is a significant difference between functional correctness (e.g., ``this code
will sort a list'') and security.  In the latter case we can only prove security
against a policy and a particular threat model---we can't make any blanket
statements about attacks that might fall outside the model.

As an example of the complexity of side-channel attacks that occur in real life,
we briefly consider Spectre \citep{Kocher19sp}.  It exploits a some features of
a modern processor, namely speculative execution, together with timing attack on
the memory cache.  They use the following sample code to illustrate their
technique:
\begin{verbatim}
  if (x < array1_size)
      y = array2[array1[x] * 4096];
\end{verbatim}
This code tries to guard access to \verb'array1' to prevent out-of-bounds
access.  In the first phase of the attack, it is given many valid values for
\verb'x', training the processor's branch prediction algorithm to believe that
this branch will usually be taken.  In the second phase, it is given and input
that is out of bounds, that is, greater than \verb'array1_size'.  Branch
prediction will assume the test turns out to be true and starts executing
\verb'y = array2[array1[x] * 4096]'.  When the processor discovers that the
condition is actually false, it will abandon all the actions taken, such as
retrieving \verb'array1[x]', multiplying it to obtain some address \verb'addr',
retrieving a value from \verb'array2[addr]', and any other registers, condition,
codes, etc.  However, one thing it does \textbf{not} undo is storing some
fetched data in the memory cache.  Note that the ``secret'' data at
\verb'array1[x]' are used as an address, so the data at this address might still
be in the cache.  Now in the third phase the attacker launches a timing attack
against the cache to determine which memory region is now in the cache, which
can reveal the underlying data at \verb'array1[x]'.

\section{Some Simple Timing Attacks}

Real timing attacks \citep{Brumley05cn} are complex, among other things due to
timing variations in processor execution.  We abstract away from these details
and it turns out that the main defenses, like ``constant-time programming'', are
ultimately the same.  Here are some timing attacks, intentionally somewhat
hypothetical (slightly outside our language) to give you the opportunity to
devise your own in Lab 2.

With $\Sigma_0 = (\mi{pin} : \ms{H})$, we write
\[
  \mb{while}\; (\mi{pin} > 0)\; \mi{pin} := \mi{pin} - 1
\]
Assuming the original value of $\mi{pin}$ (say $c$) is positive, this will take
something like $4 * c + 2$ steps to terminate.  Here we count every arithmetic
operation or comparison as 1 step, the assignment as 1 step, and processing of
$\mb{while}$ as another step.  While details may vary with the cost model, it
should be clear that it is easy to determine the secret value of $\mi{pin}$ from
the running time.

This code, however, does not represent a very useful attack, because it is too
slow, for example, if the pin has 256 bits.  But we can modify it into a
standard strawman example (omitting any details regarding possibly necessary
declassification):
\begin{tabbing}
  \quad $\mi{auth} := 1 \semi$ \\
  \quad $i := 0 \semi$ \\
  \quad $\mb{while}\; (i < \mi{len})$ \\
  \qquad $\mb{if}\; \mi{pin}[i] = \mi{guess}[i]$ \\
  \qquad $\mb{then}\; i := i+1$ \\
  \qquad $\mb{else}\; (\mi{auth} := 0 \semi i = \mi{len})$
\end{tabbing}
Here $\mi{str}[i]$ could either access the $i$th character of a string
$\mi{str}$ or the $i$th bit of the number.  $\mi{len}$ would represent either
the length of the string or the width of the pin in bits.

Let's assume there are nonnegative integers, and $\mi{len}$ is the number of
bits.  If the lowest bit of the guess is incorrect, the program will terminate
with $\mi{auth}$ almost immediately.  If it is correct, it will go around the
loop at least one more time, and therefore take longer.  That will allow us to
infer the lowest bit of the secret pin.  Our next guess will have that bit
correct and now find the next bit, etc.  With at most $2 * \mi{len}$ experiments
we should be able to determine all bits.

\section{Time-Sensitive Noninterference}

As in the last lecture, our first job will be to define the correct semantic
notion of noninterference and measure it against the examples.  We would like
that if both programs terminate from prestates that are indistinguishable for a
low-security observer, they take the same number of steps and terminate in
low-security equivalent states.  In order to express the number of steps, we
change our evaluation function to return not only a poststate, but also the
number of steps taken.  The latter will be defined shortly.
\begin{quote}
  We define $\Sigma \models \alpha\; \ms{secure}^\tm$ (program $\alpha$
  satisfies \emph{time-sensitive noninterference with respect to policy
    $\Sigma$}) \newline iff
  \begin{enumerate}[]
  \item for all $\ell$, $\omega_1$, $\omega_2$, $\nu_1$, $\nu_2$, $n_1$, and $n_2$,
  \item whenever $\Sigma \vdash \omega_1 \approx_\ell \omega_2$,
  \item and $\eval\; \omega_1\; \alpha = (n_1, \nu_1)$,
  \item and $\eval\; \omega_2\; \alpha = (n_2, \nu_2)$,
  \item then $\Sigma \vdash \nu_1 \approx_\ell \nu_2$ and $n_1 = n_2$.
  \end{enumerate}
\end{quote}
Note that we have disregarded nontermination in this definition.  It is easy to
combine it with the requirement to be termination-sensitive, but we prefer to
separate these concerns for now.

Now our first example (with $\Sigma_0 = (\mi{pin} : \ms{H})$) clearly does not
satisfy this definition, which is what we would hope.
\[
  \alpha_0 = (\mb{while}\; (\mi{pin} > 0)\; \mi{pin} := \mi{pin} - 1)
\]
A simple counterexample is
\[
  \begin{array}{l@{\hspace{3em}}l}
    \omega_1 = (\mi{pin} \mapsto 0) & \eval\; \omega_1\; \alpha_0 = (2, \mi{pin} \mapsto 0) \\
    \omega_2 = (\mi{pin} \mapsto 1) & \eval\; \omega_2\; \alpha_0 = (6, \mi{pin} \mapsto 0)
  \end{array}
\]
We have $\Sigma_0 \vdash \omega_1 \approx_{\ms{L}} \omega_2$ (since
$\mi{pin} : \ms{H}$), and the poststates are indistinguishable (not only for a
low-security observer, but in general), but the steps $n_1 = 2 \neq 6 = n_2$.

The second example is subject to a similar analysis, but we'd need to be careful
about declassification which we would like to avoid.

\section{Evaluation with Time}

Next we should define evaluation with step-counting. This is quite
straightforward, just a bit tedious.  You can find the summary of this exercise
in \autoref{fig:timed-eval}.  We have to recognize that this form of timing
isn't realistic and hope that to some extent the design of the information flow
type system can make up for this lack of realism.

One point about the Boolean operations: we do not want their running times to be
data-dependent.  This means that evaluating $\top \land P$ and $\bot \land P$
should take the same amount of time.  In other words, the Boolean operations
should not be ``short-circuiting''.  Recall that in a previous lecture we
determined that it didn't matter whether they were short-circuiting or not since
all expressions and formulas are safe and terminating.  In this context it does
matter, and steps may need to be taken to ensure data independence.

Timing considerations raise another point about abstractions.  So far we have
said that variables can hold arbitrary integers.  However, this means that the
running times of many operations like addition cannot be data-independent.  So
we assume instead that integers are implemented with a fixed range, like 64 bits
or 256 bits and that arithmetic is modular.  This creates a number of
complications when reasoning about the correctness of code, but it actually
makes reasoning about certain security properties (especially those that are
timing-sensitive) both more realistic and simpler.

\begin{figure}[ht]
  \[
    \begin{array}{lcll}
      \evalz\; \omega\; c & = & (0, c) \\
      \evalz\; \omega\; x & = & (1, \omega(x)) \\
      \evalz\; \omega\; (e_1 + e_2) & = & (n_1 + n_2 + 1, c_1 + c_2) \\
                          & & \mbox{where $\evalz\; \omega\; e_1 = (n_1, c_1)$} \\
                          & & \mbox{and $\evalz\; \omega\; e_2 = (n_2, c_2)$}
      \\[1em]
      \evalb\; \omega\; (e_1 \leq e_2) & = & (n_1 + n_2 + 1, c_1 \leq c_2) \\
                          & & \mbox{where $\evalz\; \omega\; e_1 = (n_1, c_1)$} \\
                          & & \mbox{and $\evalz\; \omega\; e_2 = (n_2, c_2)$} \\
      \evalb\; \omega\; (\top) & = & (0, \top) \\
      \evalb\; \omega\; (\bot) & = & (0, \bot) \\
      \evalb\; \omega\; (P \land Q) & = & (n_1 + n_2 + 1, b_1 \land b_2) \\
                          & & \mbox{where $\evalb\; \omega\; P = (n_1, b_1)$} \\
                          & & \mbox{and $\land \evalb\; \omega\; Q = (n_2, b_2)$}
      \\[1em]
      \eval\; \omega\; (x := e) & = & (n + 1, \omega[x \mapsto c]) \\
                          & & \mbox{where $\evalz\; \omega\; e = (n, c)$}
      \\[1ex]
      \eval\; \omega\; (\alpha \semi \beta) & = & (n_1 + n_2 + 1, \nu) \\
                          & & \mbox{where $\eval\; \omega\; \alpha = (n_1, \mu)$} \\
                          & & \mbox{and $\eval\; \mu\; \beta = (n+1, \nu)$}
      \\[1ex]
      \eval\; \omega\; (\mb{skip}) & = & (1, \omega) \\
      \\[1ex]
      \eval\; \omega\; (\mb{if}\; P\; \mb{then}\; \alpha\; \mb{else}\; \beta)
                          & = & (k + n + 1, \nu) \\
                          & & \mbox{where $\eval\; \omega\; P = (k, \top)$ and $\eval\; \omega\; \alpha = (n, \nu)$} \\
                          & & \mbox{or $\eval\; \omega\; P = (k, \bot)$ and $\eval\; \omega\; \beta = (n, \nu)$}
      \\[1ex]
      \eval\; \omega\; (\mb{while}\; P\; \alpha) & = & (k + n + 1, \nu) \\
                          & & \mbox{where $\eval\; \omega\; P = (k, \bot)$ and $n = 0$ and $\nu = \omega$} \\
                          & & \mbox{or $\eval\; \omega\; P = (k, \top)$
                              and $\eval\; \omega\; (\alpha \semi \mb{while}\; P\; \alpha) = (n, \nu)$}
    \end{array}
  \]
  \caption{Timed Evaluation}
  \label{fig:timed-eval}
\end{figure}

\section{A Type System for Constant-Time Computation}

Next comes the task to develop a type system that enforces time-sensitive
noninterference, which is usually called ``\emph{constant-time}''.  However, it
isn't actually constant time, it is just that the running time can only depend
on low-security inputs.  In order to formalize that, we once again proceed
construct by construct.  We write $\Sigma \vdash \alpha\; \ms{secure}^\tm$ for
time-sensitive security with respect to signature $\Sigma$.

\paragraph{Assignment.}  For now, there doesn't seem to be any reason to change
our rule for assignment.  The point is that the time for the evaluation of an
expression $e$ is some $n$, regardless of the state $\omega$.  So it just once
again comes down to the basic information flow properties.
\begin{rules}
  \infer[\mb{:=}F^\tm, \mbox{preliminary version}]
    {\Sigma \vdash x := e\; \ms{secure}^\tm}
    {\Sigma \vdash e : \ell
      & \ell' = \Sigma(\mi{pc}) \sqcup \ell
      & \ell' \sqsubseteq \Sigma(x)}
\end{rules}
See the end of this section for further thoughts on assignment.

\paragraph{Sequential Composition and $\mb{skip}$.}  If $\alpha$ and $\beta$ are
indistinguishable even via a timing channel, then their composition should not
be, either.  The proof is straightforward and follows familiar patterns, so we
omit it here.  $\mb{skip}$ is trivial, as usual.
\begin{rules}
  \infer[\mb{\semi}F^\tm]
  {\Sigma \vdash \alpha \semi \beta\; \ms{secure}^\tm}
  {\Sigma \vdash \alpha\; \ms{secure}^\tm
    & \Sigma \vdash \beta\; \ms{secure}^\tm}
  \hspace{3em}
  \infer[\mb{skip}F^\tm]
  {\Sigma \vdash \mb{skip}\; \ms{secure}^\tm}
  {\mathstrut}
\end{rules}

\paragraph{Conditionals.}  The first instinct might be that in a conditional,
both branches need to take the same amount of time.  Under such a discipline,
we could rewrite our motivating example as follows:
\begin{tabbing}
  \quad $\mi{auth} := 1 \semi$ \\
  \quad $i := 0 \semi$ \\
  \quad $\mb{while}\; (i < \mi{len})$ \\
  \qquad $\mb{if}\; \mi{pin}[i] = \mi{guess}[i]$ \\
  \qquad $\mb{then}\; \mi{auth} := \mi{auth}$ \\
  \qquad $\mb{else}\; \mi{auth} := 0 \semi$ \\
  \qquad $i := i+1$
\end{tabbing}
The good part here is that we go around the loop the same number of times,
regardless whether the guess is correct or not.  The bad part is that
even though the two branches, abstractly, take the same amount of time
that is unlikely to be case in practice.  Any reasonable compiler would
optimize $\mi{auth} := \mi{auth}$ into a no-op and the loop would leak
some timing information.

The problem actually goes deeper: in order to have a \emph{type system} for
time-sensitive information flow that allows such a program, the type system
would have to capture exactly how long an expression or command takes to
evaluate.  Not only would this be inaccurate with respect to the actually
running code, it would also require our type-checker to perform arithmetic
reasoning.  This would just be too brittle a design.

Instead, we take a more drastic step: for any conditional, we require that the
formula $P$ depends only on low-security variables.  This rules out our sample
program above, because the comparison $\mi{pin}[i] = \mi{guess}[i]$ depends on
$\mi{pin}$ which is of high security.

Intuitively, we are repeating the solution for termination-sensitive information
flow from the last lecture, but with conditionals instead of loops.  Here is the
first formulation: we just require $\ell'$ to be $\bot$, the least element of
the security lattice.
\begin{rules}
  \infer[\mb{if}F^\tm]
  {\Sigma \vdash \mb{if}\; P\; \mb{then}\; \alpha\; \mb{else}\; \beta\; \ms{secure}^\tm}
  {\Sigma \vdash P : \ell
    & \bot = \ell' = \Sigma(\mi{pc}) \sqcup \ell
    & \Sigma' = \Sigma[\mi{pc} \mapsto \ell']
    & \Sigma' \vdash \alpha\; \ms{secure}^\tm
    & \Sigma' \vdash \beta\; \ms{secure}^\tm}
\end{rules}
We notice that this requires $\ell = \Sigma(\mi{pc}) = \bot$.  And if
$\ell' = \bot = \ell$, there is no need to update the $\mi{pc}$, so we
can use $\Sigma' = \Sigma$.
\begin{rules}
  \infer[\mb{if}F^\tm]
  {\Sigma \vdash \mb{if}\; P\; \mb{then}\; \alpha\; \mb{else}\; \beta\; \ms{secure}^\tm}
  {\Sigma \vdash P : \bot
    & \Sigma \vdash \alpha\; \ms{secure}^\tm
    & \Sigma \vdash \beta\; \ms{secure}^\tm}
\end{rules}
At this point we realize we don't need $\mi{pc}$ at all, since we only
introduced it to track the implicit flow from the condition into the branches of
an if-then-else.

With this simplification, we have drastically reduced the range of programs that
are considered secure.  So how do we rewrite our example?  We have to ``inline''
the conditional as an ordinary operation.  Note that the number of bit (or
maximal length of the string) would have to be a low-security variable.
\begin{tabbing}
  \quad $\mi{auth} := 1 \semi$ \\
  \quad $i := 0 \semi$ \\
  \quad $\mb{while}\; (i < \mi{len})$ \\
  \qquad $\mi{auth} := \mi{auth} \land (\mi{pin}[i] = \mi{guess}[i]) \semi$ \\
  \qquad $i := i+1$
\end{tabbing}
This is just outside the range of what our language permits in that equality
produces a Boolean, and conjunction takes two Booleans, but here they are
integers.  This could be solved any number of ways.  Perhaps the simplest is
that the Boolean operators are overloaded to also work on integers, for example,
with $\bot$ represented by $0$ and $\top$ by $1$, or by anything nonzero.

Also, we go back and think of $\mi{pin}$ and $\mi{guess}$ as strings rather
than nonnegative integers, this code still makes sense and is not subject
to timing attacks.

Let's prove the soundness of this rule.  We set up:
\begin{tabbing}
  $\Sigma \models P : \bot$ \` (1, first premise) \\
  $\Sigma \models \alpha\; \ms{secure}^\tm$ \` (2, second premise) \\
  $\Sigma \models \beta\; \ms{secure}^\tm$ \` (3, third premise) \\
  $\Sigma \vdash \omega_1 \approx_\ell \omega_2$ \` (4, assumption) \\
  $\eval\; \omega_1\; (\mb{if}\; P\; \mb{then}\; \alpha\; \mb{else}\; \beta) = (n_1, \nu_1)$ \` (5, assumption) \\
  $\eval\; \omega_2\; (\mb{if}\; P\; \mb{then}\; \alpha\; \mb{else}\; \beta) = (n_2, \nu_2)$ \` (6, assumption) \\
  $\ldots$ \\
  $\Sigma \vdash \nu_1 \approx_\ell \nu_2$ and \` (to show) \\
  $n_1 = n_2$ \` (to show)
\end{tabbing}
Because $\Sigma \vdash \omega_1 \approx_\ell \omega_2$ and
$\bot \sqsubseteq \ell$, we have that
$\eval\; \omega_1\; P = \eval\; \omega_2\; P = (k, b)$ for some nonnegative time
$k$ and Boolean $b$.  We consider the case where $b = \top$; the other is
symmetric.  Then
\[
  \eval\; \omega_1\; (\mb{if}\; P\; \mb{then}\; \alpha\; \mb{else}\; \beta) = (k + m_1
  + 1, \nu_1)
\]
where $\eval\; \omega_1\; \alpha = (m_1, \nu_1)$ and $n_1 = k + m_1 + 1$.

Also
\[
  \eval\; \omega_2\; (\mb{if}\; P\; \mb{then}\; \alpha\; \mb{else}\; \beta) = (k + m_2
  + 1, \nu_2)
\]
where $\eval\; \omega_2\; \alpha = (m_2, \nu_2)$ and $n_2 = k + m_2 + 1$.

Since $\alpha$ is secure (that is, $\Sigma \models \alpha\; \ms{secure}^\tm$,
which comes from the second premise) we get
$\Sigma \vdash \nu_1 \approx_\ell \nu_2$ (which is one of the two properties we
needed to prove) and $m_1 = m_2$.

Therefore, also $n_1 = k + m_1 + 1 = k + m_2 + 1 = n_2$, which is the second
property we needed to prove.

\paragraph{Loops.}  For the same reason as conditionals, we require
the loop guard to be of low security.
\begin{rules}
  \infer[\mb{while}F^\tm]
  {\Sigma \vdash \mb{while}\; P\; \alpha\; \ms{secure}^\tm}
  {\Sigma \vdash P : \bot
    & \Sigma \vdash \alpha\; \ms{secure}^\tm}
\end{rules}

\paragraph{Tests.} If the test fails and the program aborts we have a different
kind of channel.  Again, the simplest condition that avoid information flow here
(even if outside of the definition) is to require the test to be of low
security.
\begin{rules}
  \infer[\mb{test}F]
  {\Sigma \vdash \mb{test}\; P\; \ms{secure}^\tm}
  {\Sigma \vdash P : \bot}
\end{rules}

\paragraph{Assignment Revisited.}  Since we eliminated the ghost variable $\mi{pc}$,
we simplify the rule for assignment.
\begin{rules}
  \infer[\mb{:=}F^\tm]
    {\Sigma \vdash x := e\; \ms{secure}^\tm}
    {\Sigma \vdash e : \ell
      & \ell \sqsubseteq \Sigma(x)}
\end{rules}

\begin{figure}[ht]
  \begin{rules}
    \infer[\mb{:=}F^\tm]
    {\Sigma \vdash x := e\; \ms{secure}^\tm}
    {\Sigma \vdash e : \ell
      & \ell \sqsubseteq \Sigma(x)}
    \\[1em]
    \infer[\mb{\semi}F^\tm]
    {\Sigma \vdash \alpha \semi \beta\; \ms{secure}^\tm}
    {\Sigma \vdash \alpha\; \ms{secure}^\tm
      & \Sigma \vdash \beta\; \ms{secure}^\tm}
    \hspace{3em}
    \infer[\mb{skip}F^\tm]
    {\Sigma \vdash \mb{skip}\; \ms{secure}^\tm}
    {\mathstrut}
    \\[1em]
    \infer[\mb{if}F^\tm]
    {\Sigma \vdash \mb{if}\; P\; \mb{then}\; \alpha\; \mb{else}\; \beta\; \ms{secure}^\tm}
    {\Sigma \vdash P : \bot
      & \Sigma \vdash \alpha\; \ms{secure}^\tm
      & \Sigma \vdash \beta\; \ms{secure}^\tm}
    \hspace{3em}
    \infer[\mb{while}F^\tm]
    {\Sigma \vdash \mb{while}\; P\; \alpha\; \ms{secure}^\tm}
    {\Sigma \vdash P : \bot
      & \Sigma \vdash \alpha\; \ms{secure}^\tm}
    \\[1em]
    \infer[\mb{test}F]
    {\Sigma \vdash \mb{test}\; P\; \ms{secure}^\tm}
    {\Sigma \vdash P : \bot}
  \end{rules}
  \caption{Timing Sensitive Information Flow Typing}
  \label{fig:timing-sensitive}
\end{figure}

A summary of the rules can be found in \autoref{fig:timing-sensitive}.  It is
remarkably simple because it is not afraid to rule out many programs.  The
soundness of all the rules taken together gives us the soundness of the whole
type system.

\begin{theorem}[Soundness of timing-sensitive information flow typing]
  \mbox{}\newline
  If $\Sigma \vdash \alpha\; \ms{secure}^\tm$ then
  $\Sigma \models \alpha\; \ms{secure}^\tm$
\end{theorem}
\begin{proof}
  All the rules are sound in the sense that they preserve semantic validity of
  the premises.  By a trivial induction over the derivation of
  $\Sigma \vdash \alpha\; \ms{secure}^\tm$ we obtain
  $\Sigma \models \alpha\; \ms{secure}^\tm$.
\end{proof}

Comparing the type system against the one from the last lecture for
termination-sensitive noninterference, we see that it actually enforces that as
well.  One can update the definition of noninterference and statement of the
above theorem to account for that.

\section{Randomization}

Another defense against timing attacks is to introduce randomization into the
program.  Even though an observer might still be able to observe timing, this
information will then be insufficient to determine the secret.  In our example
(using the interpretation of pins and guesses as large integers) we could write
something like:
\begin{tabbing}
  \quad $r := \ms{random}() \semi$ \\
  \quad $\mi{pin} := \mi{pin} + r \semi$ \\
  \quad $\mi{guess} := \mi{guess} + r \semi$ \\
  \quad $\mi{auth} := 1 \semi$ \\
  \quad $i := 0 \semi$ \\
  \quad $\mb{while}\; (i < \mi{len})$ \\
  \qquad $\mb{if}\; \mi{pin}[i] = \mi{guess}[i]$ \\
  \qquad $\mb{then}\; i := i+1$ \\
  \qquad $\mb{else}\; (\mi{auth} := 0 \semi i = \mi{len})$
\end{tabbing}
Every time this program is run, it is like to use a different random number,
essentially making it impossible to draw useful conclusions from the running
times.  Some care must be taken for the randomization and its result to be
sufficiently uniform.

\section{Practice Problems}

These problems are for self-study and are not to be handed in.

\subsection{Timing Analysis [3 parts]}

Consider the two-level security lattice with $\ms{L} \sqsubset \ms{H}$ and
the following security policy:
\[
  \Sigma = (\mi{pin} : \ms{H},\; \mi{guess} : \ms{L},\; \mi{auth} : \ms{L},\; i : \ms{L},\; n : \ms{L})
\]
For each of the following programs, determine whether it satisfies
time-sensitive noninterference ($\Sigma \models \alpha\; \ms{secure}^\tm$) and
whether it is derivable in the timing-sensitive type system
($\Sigma \vdash \alpha\; \ms{secure}^\tm$).  Justify each answer with either
a derivation, a brief argument, or a counterexample as appropriate.  Assume
that the timed evaluation semantics from \autoref{fig:timed-eval} are used.

\begin{task}\rm
  \label{task:p1a}
  $\alpha_1 = (n := \mi{pin} + \mi{guess} \semi \mi{auth} := \mi{guess})$
\end{task}

\begin{task}\rm
  \label{task:p1b}
  $\alpha_2 = (\mb{if}\; \mi{guess} > 0\; \mb{then}\; \mi{auth} := 1\;
  \mb{else}\; (\mi{auth} := 0 \semi \mi{auth} := \mi{auth}))$
\end{task}

\begin{task}\rm
  \label{task:p1c}
  $\alpha_3 = (\mb{while}\; (i < n)\;
  (\mi{auth} := \mi{auth} \land (\mi{pin} + i) \semi i := i + 1))$
\end{task}

\subsection{Constant-Time Conditional [4 parts]}

In constant-time programming, conditionals that depend on secret data are
typically replaced by \emph{conditional moves} or \emph{select} operations.
We add a new construct $\mb{csel}(P,\; e_1,\; e_2)$ to our expression language.
Its semantics are: $\evalz\; \omega\; \mb{csel}(P, e_1, e_2) = (n_P + n_1 + n_2 + 1, c)$
where $\evalb\; \omega\; P = (n_P, b)$, $\evalz\; \omega\; e_1 = (n_1, c_1)$,
$\evalz\; \omega\; e_2 = (n_2, c_2)$, and
$c = c_1$ if $b = \top$, $c = c_2$ if $b = \bot$.

The key feature of $\mb{csel}$ is that it \emph{always evaluates both
  expressions}, so the time cost $n_P + n_1 + n_2 + 1$ does not depend on which
branch is taken.

\begin{task}\rm
  \label{task:p2a}
  What is the security level of $\mb{csel}(P, e_1, e_2)$?  That is, propose a
  rule of the form
  \[
    \infer{\Sigma \vdash \mb{csel}(P, e_1, e_2) : \ell}
    {\ldots}
  \]
  Your rule should be as permissive as possible.  In particular, consider
  whether $P$ needs to be of level $\bot$ as it does for $\mb{if}$ in the
  timing-sensitive type system.
\end{task}

\begin{task}\rm
  \label{task:p2b}
  Prove that your rule for $\mb{csel}$ is sound for time-sensitive
  noninterference.  That is, show that if $\Sigma \vdash \mb{csel}(P, e_1, e_2)
  : \ell$ and $\Sigma \vdash \omega_1 \approx_\ell \omega_2$, then
  $\evalz\; \omega_1\; \mb{csel}(P, e_1, e_2)$ and
  $\evalz\; \omega_2\; \mb{csel}(P, e_1, e_2)$ yield the same time cost and
  the same value.
\end{task}

\begin{task}\rm
  \label{task:p2c}
  Consider the following program under the policy
  $\Sigma' = (\mi{pin} : \ms{H},\; \mi{guess} : \ms{L},\; r : \ms{L})$:
  \[
    r := \mb{csel}(\mi{pin} > 0,\; \mi{guess} + 1,\; \mi{guess})
  \]
  Does this program satisfy $\Sigma' \models \alpha\; \ms{secure}^\tm$?  Is it
  derivable as $\Sigma' \vdash \alpha\; \ms{secure}^\tm$ using the
  timing-sensitive type system extended with your rule from \autoref{task:p2a}?
  Justify your answers.
\end{task}

\begin{task}\rm
  \label{task:p2d}
  Give an example of a program that uses $\mb{csel}$ and satisfies time-sensitive
  noninterference semantically ($\Sigma \models \alpha\; \ms{secure}^\tm$), but
  which would \emph{not} satisfy it if $\mb{csel}(P, e_1, e_2)$ were replaced
  by $\mb{if}\; P\; \mb{then}\; x := e_1\; \mb{else}\; x := e_2$ (under the same
  policy).  Briefly explain the difference.
\end{task}

\bibliographystyle{plainnat}
\bibliography{bibliography}

\end{document}
